\section{Physics Layer II: Bifurcations, Thresholds, and Critical Behaviour}

\begin{definition}[Bifurcation-style change]
A \textbf{bifurcation-style change} occurs when a parameter shift changes the
qualitative structure of the dynamics rather than merely the numerical speed of
movement.
\end{definition}

Bifurcation language matters because many interventions are valuable precisely
when they are small in effort but large in structural consequence. A new alarm,
a changed route, a blocked app, or a pre-written task cue may look minor in
energy cost yet large in outcome because it has moved a control parameter past a
threshold.

\begin{example}[Commute-to-gym routing]
If a worker goes home first, the evening may fall into a rest basin. If the
route is changed so that the commute passes through the gym entrance, the
system may never reconstruct the old basin at all. The intervention is small in
description but large in consequence because it changes when and where the
decision is made. That is exactly the sort of regime-sensitive reasoning
bifurcation language is meant to capture.
\end{example}
